\documentclass{beamer}

\usepackage[italian]{babel}
\usepackage[utf8]{inputenc}
\usepackage{listings}
\usepackage{xcolor}

\usetheme{Madrid}

\title{Trigger nei Database Relazionali}
\author{Prof. Massimo Fedeli}
\date{\today}

\lstset{
	language=SQL,
	basicstyle=\ttfamily\small,
	keywordstyle=\color{blue},
	commentstyle=\color{gray},
	stringstyle=\color{red},
	showstringspaces=false,
	breaklines=true
}

\begin{document}
	
	\frame{\titlepage}
	
	%---------------------------------
	\begin{frame}{Definizione di Trigger}
		Un \textbf{trigger} è un programma eseguito automaticamente in risposta a un evento su una tabella.
		
		\bigskip
		
		Eventi tipici:
		\begin{itemize}
			\item INSERT
			\item UPDATE
			\item DELETE
		\end{itemize}
		
		\bigskip
		
		I trigger sono parte integrante della logica lato database.
	\end{frame}
	
	%---------------------------------
	\begin{frame}{Tipologie di Trigger}
		Lo standard SQL definisce due categorie:
		
		\begin{itemize}
			\item \textbf{Trigger a livello di riga}
			\item \textbf{Trigger a livello di istruzione}
		\end{itemize}
		
		\bigskip
		
		Differenza fondamentale:
		\begin{itemize}
			\item Row-level: eseguito per ogni riga
			\item Statement-level: eseguito una sola volta
		\end{itemize}
	\end{frame}
	
	%---------------------------------
	\begin{frame}{Trigger a Livello di Riga}
		Viene eseguito per ogni riga modificata.
		
		\bigskip
		
		Esempio:
		\begin{itemize}
			\item UPDATE su 100 righe $\rightarrow$ 100 esecuzioni
		\end{itemize}
		
		\bigskip
		
		Utilizzo tipico:
		\begin{itemize}
			\item validazione dati
			\item logging dettagliato
		\end{itemize}
	\end{frame}
	
	%---------------------------------
	\begin{frame}{Trigger a Livello di Istruzione}
		Viene eseguito una sola volta per ogni istruzione SQL.
		
		\bigskip
		
		Esempio:
		\begin{itemize}
			\item UPDATE su 100 righe $\rightarrow$ 1 esecuzione
		\end{itemize}
		
		\bigskip
		
		Nota:
		\begin{itemize}
			\item MySQL \textbf{non supporta} questo tipo
		\end{itemize}
	\end{frame}
	
	%---------------------------------
	\begin{frame}{Supporto nei DBMS}
		MySQL:
		\begin{itemize}
			\item Supporta solo trigger \textbf{FOR EACH ROW}
			\item Non supporta trigger statement-level
		\end{itemize}
		
		\bigskip
		
		Altri DBMS (PostgreSQL, Oracle):
		\begin{itemize}
			\item Supportano entrambe le modalità
		\end{itemize}
	\end{frame}
	
	%---------------------------------
	\begin{frame}{Vantaggi dei Trigger}
		\begin{itemize}
			\item Controllo dell'integrità dei dati
			\item Automazione di operazioni
			\item Logging delle modifiche
			\item Riduzione logica lato applicazione
		\end{itemize}
	\end{frame}
	
	%---------------------------------
	\begin{frame}{Quando Usarli}
		Uso appropriato:
		\begin{itemize}
			\item auditing
			\item validazione automatica
			\item aggiornamento tabelle derivate
		\end{itemize}
		
		\bigskip
		
		Uso improprio:
		\begin{itemize}
			\item logiche troppo complesse
			\item dipendenze difficili da tracciare
		\end{itemize}
	\end{frame}
	
	%---------------------------------
	\begin{frame}[fragile]{Sintassi CREATE TRIGGER}
		\begin{lstlisting}
CREATE TRIGGER trigger_name
BEFORE | AFTER INSERT | UPDATE | DELETE
ON table_name
FOR EACH ROW
BEGIN
  -- istruzioni SQL
END;
		\end{lstlisting}
	\end{frame}
	
	%---------------------------------
	\begin{frame}{Componenti del Trigger}
		\begin{itemize}
			\item Nome del trigger
			\item Evento (INSERT, UPDATE, DELETE)
			\item Timing (BEFORE, AFTER)
			\item Tabella associata
			\item Corpo del trigger
		\end{itemize}
	\end{frame}
	
	%---------------------------------
	\begin{frame}[fragile]{Esempio: BEFORE INSERT}
		\begin{lstlisting}
CREATE TRIGGER before_insert_user
BEFORE INSERT ON users
FOR EACH ROW
BEGIN
  SET NEW.created_at = NOW();
END;
		\end{lstlisting}
		
		\bigskip
		
		Imposta automaticamente il timestamp.
	\end{frame}
	
	%---------------------------------
	\begin{frame}[fragile]{Esempio: AFTER INSERT}
		\begin{lstlisting}
CREATE TRIGGER after_insert_order
AFTER INSERT ON orders
FOR EACH ROW
BEGIN
  INSERT INTO logs(message)
  VALUES ('Nuovo ordine inserito');
END;
		\end{lstlisting}
	\end{frame}
	
	%---------------------------------
	\begin{frame}[fragile]{Esempio: BEFORE UPDATE}
		\begin{lstlisting}
CREATE TRIGGER before_update_salary
BEFORE UPDATE ON employees
FOR EACH ROW
BEGIN
  IF NEW.salary < 0 THEN
    SIGNAL SQLSTATE '45000';
  END IF;
END;
		\end{lstlisting}
	\end{frame}
	
	%---------------------------------
	\begin{frame}[fragile]{Esempio: AFTER UPDATE}
		\begin{lstlisting}
CREATE TRIGGER after_update_employee
AFTER UPDATE ON employees
FOR EACH ROW
BEGIN
  INSERT INTO audit_log(old_salary, new_salary)
  VALUES (OLD.salary, NEW.salary);
END;
		\end{lstlisting}
	\end{frame}
	
	%---------------------------------
	\begin{frame}[fragile]{Esempio: BEFORE DELETE}
		\begin{lstlisting}
CREATE TRIGGER before_delete_user
BEFORE DELETE ON users
FOR EACH ROW
BEGIN
  INSERT INTO archive_users
  VALUES (OLD.id, OLD.name);
END;
		\end{lstlisting}
	\end{frame}
	
	%---------------------------------
	\begin{frame}[fragile]{Esempio: AFTER DELETE}
		\begin{lstlisting}
CREATE TRIGGER after_delete_order
AFTER DELETE ON orders
FOR EACH ROW
BEGIN
  INSERT INTO logs(message)
  VALUES ('Ordine eliminato');
END;
		\end{lstlisting}
	\end{frame}
	
	%---------------------------------
	\begin{frame}{Uso di OLD e NEW}
		\begin{itemize}
			\item \texttt{NEW}: valori dopo la modifica
			\item \texttt{OLD}: valori prima della modifica
		\end{itemize}
		
		\bigskip
		
		Disponibilità:
		\begin{itemize}
			\item INSERT $\rightarrow$ solo \texttt{NEW}
			\item DELETE $\rightarrow$ solo \texttt{OLD}
			\item UPDATE $\rightarrow$ \texttt{OLD} e \texttt{NEW}
		\end{itemize}
	\end{frame}
	
	%---------------------------------
	\begin{frame}[fragile]{Eliminazione di un Trigger}
		\begin{lstlisting}
DROP TRIGGER IF EXISTS trigger_name;
		\end{lstlisting}
		
		\bigskip
		
		Note:
		\begin{itemize}
			\item \texttt{IF EXISTS} evita errori
			\item servono privilegi adeguati
		\end{itemize}
	\end{frame}
	
	%---------------------------------
	\begin{frame}{Comportamento su DROP TABLE}
		Se una tabella viene eliminata:
		
		\begin{itemize}
			\item tutti i trigger associati vengono eliminati automaticamente
		\end{itemize}
	\end{frame}
	
	%---------------------------------
	\begin{frame}{Trigger Multipli}
		MySQL 8.0 consente:
		
		\begin{itemize}
			\item più trigger per stesso evento
			\item stesso timing (es. AFTER INSERT)
		\end{itemize}
		
		\bigskip
		
		Ordine di esecuzione:
		\begin{itemize}
			\item definito dal DBMS
		\end{itemize}
	\end{frame}
	
	%---------------------------------
	\begin{frame}{Caso di Studio}
		Scenario:
		\begin{itemize}
			\item tabella ordini
			\item aggiornamento automatico statistiche
		\end{itemize}
		
		\bigskip
		
		Soluzione:
		\begin{itemize}
			\item trigger AFTER INSERT
		\end{itemize}
	\end{frame}
	
	%---------------------------------
	\begin{frame}[fragile]{Esempio Completo}
		\begin{lstlisting}
CREATE TRIGGER update_total
AFTER INSERT ON orders
FOR EACH ROW
BEGIN
  UPDATE stats
  SET total_orders = total_orders + 1;
END;
		\end{lstlisting}
	\end{frame}
	
	%---------------------------------
	\begin{frame}{Conclusioni}
		\begin{itemize}
			\item Strumento potente ma da usare con cautela
			\item Ideale per automazioni lato DB
			\item Può complicare il debugging
		\end{itemize}
		
		\bigskip
		
		Regola pratica:
		\begin{itemize}
			\item usare solo quando realmente necessario
		\end{itemize}
	\end{frame}
	
\end{document}
